\subsection{Useful Modules}

\begin{itemize}
    \item \textbf{math.gcd / math.lcm: } \texttt{math.gcd(a, b)} and \texttt{math.lcm(a, b)} (Python 3.9+). \texttt{math.gcd} accepts multiple args in 3.9+: \texttt{math.gcd(*lst)}.
    \item \textbf{math.isqrt: } Integer square root in O(1), no float precision issues. Safer than \texttt{int(math.sqrt(n))}.
    \item \textbf{math.comb / math.perm: } \texttt{math.comb(n, k)} = $\binom{n}{k}$ and \texttt{math.perm(n, k)} = $P(n,k)$, exact integers, no precomputation needed.
    \item \textbf{collections.Counter: } Counts element frequencies. Supports arithmetic (\texttt{+}, \texttt{-}, \texttt{\&}, \texttt{|}). \texttt{most\_common(k)} returns top \texttt{k} elements.
    \item \textbf{collections.defaultdict: } Like a dict but returns a default value for missing keys. \texttt{defaultdict(list)}, \texttt{defaultdict(int)}, \texttt{defaultdict(set)}.
    \item \textbf{collections.deque: } Double-ended queue with O(1) \texttt{appendleft} / \texttt{popleft}. Use instead of \texttt{list} for BFS queues. Supports \texttt{maxlen} for sliding windows.
    \item \textbf{heapq: } Min-heap by default. For max-heap, negate values: push \texttt{-x}, pop and negate result. \texttt{heapq.nlargest(k, lst)} and \texttt{nsmallest(k, lst)} in O(n log k).
    \item \textbf{bisect: } \texttt{bisect\_left(a, x)} / \texttt{bisect\_right(a, x)} return insertion points in a sorted list in O(log n). \texttt{insort} inserts maintaining order.
    \item \textbf{itertools.permutations / combinations: } \texttt{permutations(lst, r)} and \texttt{combinations(lst, r)} generate lazily. \texttt{combinations\_with\_replacement} allows repeats.
    \item \textbf{itertools.product: } \texttt{product(A, B)} = Cartesian product. \texttt{product(A, repeat=k)} = $A^k$. Replaces k nested loops.
    \item \textbf{itertools.accumulate: } \texttt{accumulate(lst)} builds prefix sums. Pass \texttt{func} for other operations (e.g., \texttt{max}, \texttt{operator.mul}).
    \item \textbf{fractions.Fraction: } Exact rational arithmetic. \texttt{Fraction(1, 3) + Fraction(1, 6) = Fraction(1, 2)}. Avoids floating point errors.
    \item \textbf{string module: } \texttt{string.ascii\_lowercase}, \texttt{string.ascii\_uppercase}, \texttt{string.digits} give ready-made character sets.
\end{itemize}

\begin{lstlisting}[language=Python]
from collections import Counter, defaultdict, deque
from itertools import permutations, combinations, product, accumulate
import heapq, math, bisect

# --- Counter ---
freq = Counter("abracadabra")
print(freq.most_common(2))    # [('a',5),('b',2)]
# Check anagram
Counter("listen") == Counter("silent")  # True

# --- defaultdict ---
graph = defaultdict(list)
graph[0].append(1)            # no KeyError

# --- deque (BFS) ---
q = deque([start])
while q:
    v = q.popleft()           # O(1)

# --- heapq (max-heap trick) ---
heap = []
heapq.heappush(heap, -5)
top = -heapq.heappop(heap)   # 5

# --- bisect ---
a = [1, 3, 5, 7]
i = bisect.bisect_left(a, 4) # 2  (first pos >= 4)

# --- prefix sums ---
pre = list(accumulate([1,2,3,4]))  # [1,3,6,10]

# --- itertools.product (k nested loops) ---
for x, y, z in product(range(3), repeat=3):
    pass  # 27 combinations

# --- math ---
print(math.isqrt(10))         # 3 (safe integer sqrt)
print(math.comb(10, 3))       # 120

# --- Fraction ---
from fractions import Fraction
print(Fraction(1,3) + Fraction(1,6))  # 1/2
\end{lstlisting}